\section{Uniformity and the growth alternatives}
\label{sec:preliminaries}

We record the reductions needed to use finite arguments uniformly on
the infinite graph. Graph distance always counts edges. Unless
specified otherwise, balls are induced subgraphs as well as vertex
sets. The maximum degree is denoted by $D$.

\subsection{Graph conventions and the threshold}

It suffices to prove Theorem~\ref{thm:main} for simple graphs.
Indeed, loops have no effect on clusters. If parallel edges occur,
subdivide each edge once, using a distinct new vertex for each edge.
The resulting graph is simple and remains connected and locally
finite. There are finitely many edge orbits, since there are finitely
many vertex orbits and each representative has finite degree.
Lifting the original automorphisms therefore proves quasi-transitivity
of the subdivided graph. At parameter $t$, each original edge is
fully traversable with probability $t^2$, independently of the others.
An infinite cluster in the subdivided graph contains infinitely many
original vertices, by local finiteness. Thus percolation on it occurs
exactly when the effective configuration on the original graph
percolates, and its critical probability is $\sqrt{\pc(G)}$.
The simple-graph conclusion at this parameter gives the desired
conclusion on $G$.

Henceforth the infinite graph is simple. Finite auxiliary graphs may
have parallel edges after contractions. Local finiteness and
connectedness imply countability. Finitely many vertex orbits imply $D<\infty$.
For $p<1/D$, the probability of an open simple path of length $n$
starting at $x$ is at most $D^np^n$, which tends to zero. An infinite
cluster would contain such a path for every $n$, so
\begin{equation}
\label{eq:positive-threshold}
 \pc\ge 1/D>0.
\end{equation}

The event that some infinite cluster exists is unchanged by altering
finitely many edge states: adding finitely many edges joins only
finitely many clusters, and deleting finitely many edges splits an
infinite cluster into finitely many pieces. It is therefore a tail
event and has probability zero or one. Monotonicity in $p$ shows that
our definition of $\pc$ agrees with the supremum of the parameters
at which almost surely every cluster is finite. In particular all
clusters are finite almost surely for $p<\pc$.

We use Harris's inequality: increasing functions of independent
edge states have nonnegative covariance \cite{Harris1960}.
For finitely many bits, this follows by induction from the covariance
decomposition after conditioning on one bit; the two conditional
means are themselves increasing. For bounded functions of countably
many bits, conditional expectations on finite sets of coordinates
and martingale convergence give the same conclusion. Decreasing
functions are likewise positively correlated; an increasing and a
decreasing function are negatively correlated.

If critical percolation occurs with positive probability, countability
gives a vertex $x$ with $\PP_{\pc}(|C_x|=\infty)>0$.
For a fixed path from $y$ to $x$, Harris's inequality bounds
$\PP_{\pc}(|C_y|=\infty)$ below by this positive probability times
$\pc$ to the path length. Finitely many vertex orbits then give
\begin{equation}
\label{eq:uniform-jump}
 \theta:=\min_{x\in V}\PP_{\pc}(|C_x|=\infty)>0.
\end{equation}
This uniform lower bound is the contradiction hypothesis in both
subexponential arguments.

\subsection{Ball comparisons}

\begin{lemma}
\label{lem:orbit-growth}
There is an integer $C_0$ such that every vertex orbit meets every
ball of radius $C_0$. Moreover, for nonnegative integer radii,
\begin{equation}
\label{eq:ball-comparisons}
 V_+(r)\le V_-(r+C_0),\qquad
 V_+(r+s)\le V_+(r)V_+(s).
\end{equation}
The limit $\gamma=\lim_{r\to\infty}V_+(r)^{1/r}$ exists. If
$\gamma>1$, every vertex ball has exponential growth; if $\gamma=1$,
then $\log V_+(r)=o(r)$.
\end{lemma}

\begin{proof}
Choose a representative in each of the finitely many orbits. For
each representative and each orbit, fix a path to that orbit, and
let $C_0$ bound the lengths of these finitely many paths. Translating
them proves the density assertion. To compare balls, move the center
of a ball attaining $V_+(r)$ within distance $C_0$ of a center attaining
$V_-(r+C_0)$. The translated smaller ball lies in the larger one.
Also $B(x,r+s)$ is covered by the radius-$s$ balls centered in
$B(x,r)$, proving the second inequality.

Thus $\log V_+(r)$ is subadditive. Dividing a radius into blocks of
any fixed length shows that its quotient by $r$ converges to its
infimum. This proves existence of $\gamma$. Finally
$V_+(r-C_0)\le |B(x,r)|\le V_+(r)$ for $r\ge C_0$ gives the claimed
uniform growth conclusions.
\end{proof}

We now invoke precisely the following external result.
\begin{theorem}[Hutchcroft {\cite[Theorem~1.2]{Hutchcroft2016}}]
\label{thm:exponential}
If a connected locally finite quasi-transitive graph satisfies
$\liminf_{r\to\infty}|B(x,r)|^{1/r}>1$, then it has almost surely no
infinite Bernoulli bond cluster at its critical probability.
\end{theorem}

It remains to work under $\log V_+(r)=o(r)$. Either
\begin{equation}
\label{eq:superpoly-growth}
 \frac{\log V_-(r)}{\log r}\longrightarrow\infty,
\end{equation}
or this quotient is bounded along an unbounded sequence of integer
radii. In the latter case there is a finite $d$ such that
\begin{equation}
\label{eq:polynomial-subsequence}
 V_-(r)\le r^d
 \quad\text{along that sequence.}
\end{equation}
These alternatives are exhaustive, including when the quotient
oscillates.

\subsection{Exact clusters and finite witnesses}

Conditioning on the exact vertex set of a cluster fixes its boundary
edges closed and imposes only an internal connectivity condition.
Consequently the edges induced outside that set keep their product
law. Revealing every incident edge, including internal edges,
has the same property. We will specify which of these two exposures
is being used.

The event $|C_x|\ge s$, for a fixed positive integer $s$, can be
witnessed inside $B(x,s-1)$: a connected cluster of at least $s$
vertices contains a connected set of $s$ vertices containing $x$.
Its probability is therefore continuous in $p$. A connection between
fixed finite sets is the increasing union of connections in finite
induced subgraphs. Its probability can be approximated from below
by such finite events, at the same parameter $p$. These observations
also apply uniformly to finitely many vertex or pair orbits.
